\subsection{Refinements of the allocation-existence comparisons}
\label{sec:allocation-refinements}

The comparisons below isolate complementary feasibility, dependence between
vertex choices, and positivity in small support dimensions. They concern
existence of some witness, not preservation of a prescribed witness.
All input classes are finite, nonempty and ample. Complements are taken in
the full ambient cube $Q_X=\{0,1\}^X$, before suppressing constant coordinates.
An empty complementary shore has its unique empty allocation and map.

\paragraph{The complete cube equations.}
For $v\in Q_X$ and $S\subseteq X$, let
$Q(v,S)=\{w\in Q_X:w|_{X\setminus S}=v|_{X\setminus S}\}$.
Write $\mathcal B(C)$ for the distinct positioned cubes contained in $C$,
including vertices, and
$I(C)=\{(v,S):v\in C,\ Q(v,S)\subseteq C\}$.
The full native matrix and allocation set are
\begin{align}
 M_C(B,(v,S))&=\mathbf1\{v\in B,\ \operatorname{dir}(B)\subseteq S\},
 \label{eq:refinement-native-matrix}\\
 \mathcal A(C)&=\{\mu\in\mathbb R^{I(C)}:\mu\ge0,\ M_C\mu=\mathbf1\}.
 \label{eq:refinement-allocation-set}
\end{align}
Here $\operatorname{dir}(B)$ is the free-coordinate set of $B$.
The vertex equations make each row at a root a probability distribution,
so $\mathcal A(C)$ is exactly the fractional formulation already used above.
Integral points correspond to representation maps by the inherited native
allocation dictionary. In particular $C\in\mathsf{Frac}$ means
$\mathcal A(C)\ne\varnothing$.

\begin{definition}[Complementary fractional and mixed feasibility]
\label{def:refinement-mixed}
For $D=Q_X\setminus C$, define
\[
 \begin{aligned}
 \mathsf{BFrac}&=\{C:\mathcal A(C)\ne\varnothing,
                         \ \mathcal A(D)\ne\varnothing\},\\
 \mathsf M&=\{C:C\in\mathsf{Rep},\ \mathcal A(D)\ne\varnothing\}.
 \end{aligned}
\]
The notation is retained from
\href{../containment-bfrac/sections/complementary_pairs.tex}{the complementary-allocation note}.
\end{definition}

\begin{proposition}[The mixed comparison]
\label{prop:refinement-mixed-chain}
One has
\[
 \mathsf{Ext}\subseteq\mathsf M\subseteq\mathsf{Rep},\qquad
 \mathsf M=\mathsf{Rep}\cap\mathsf{BFrac},\qquad
 \mathsf M\subseteq\mathsf{BFrac}\subseteq\mathsf{Frac}.
\]
Thus $\mathsf{Ext}=\mathsf{Rep}$ is equivalent to the conjunction of
$\mathsf{Ext}=\mathsf M$ and $\mathsf M=\mathsf{Rep}$.
\end{proposition}
\begin{proof}
An ordinary map gives an integral allocation. Maps on both shores
therefore imply the defining conditions of $\mathsf M$.
A represented original shore is already fractionally feasible, so its
mixed condition is precisely membership in $\mathsf{BFrac}$ as well.
The remaining inclusions and the equality equivalence follow by sandwiching.
\end{proof}

\paragraph{Joint choices that are valid on small cubes.}
A support selection $R$ assigns an actual support $R_v$ to every $v\in C$.
Its owner count on an actual cube $B$ is
\[
 b_R(B)=|\{v\in B:\operatorname{dir}(B)\subseteq R_v\}|.
\]
It is rooted at $q\in C$ if $R_q=\varnothing$ and
$R_v\cap\{x:v(x)\ne q(x)\}\ne\varnothing$ for every $v\ne q$.
Following a direction in this intersection gives an actual edge decreasing
Hamming distance to $q$. Hence each selection has a unique empty support
and a geodesic route from every vertex to that root.
Put $a_R(v,S)=\mathbf1\{R_v=S\}$.

\begin{definition}[Rooted face-consistent coupling classes]
\label{def:refinement-edge-coupling}
For an integer $k\ge0$, let $\Omega_k(C)$ be the rooted support selections
with $b_R(B)=1$ on every actual cube of dimension at most $k$, and put
\[
 \mathcal F_k(C)=\mathcal A(C)\cap
              \operatorname{conv}\{a_R:R\in\Omega_k(C)\},
 \qquad
 \mathsf E_k=\{C:\mathcal F_k(C)\ne\varnothing\}.
\]
The convex hull of an empty set is empty. The notation $\Omega_k$ and
$\mathcal F_k$ is retained from
\href{../containment-rep-frac/sections/face_consistent_coupling.tex}{RF's face-consistency note};
$\mathsf E_k$ names its class-level existence condition.
\end{definition}

Thus one joint law makes all $k$-dimensional and smaller budgets valid in
each sample, while its mean satisfies \emph{every} cube equation. The root
and mean may be chosen anew. In $\mathsf E_1$ every edge has exactly one
outgoing endpoint in each sample; $\mathsf E_2$ also has one owner on every
positioned square. These conditions are simultaneous, not separate local laws.

\begin{proposition}[The coupling hierarchy]
\label{prop:refinement-coupling-chain}
For every $k\ge0$,
\[
 \mathsf{Rep}\subseteq\mathsf E_{k+1}\subseteq\mathsf E_k
 \subseteq\mathsf{Frac}.
\]
The inherited common-root coupling theorem gives $\mathsf E_0=\mathsf{Frac}$.
On classes with largest actual cube dimension at most $k$,
$\mathsf E_k=\mathsf{Rep}$. Consequently
$\bigcap_{k\ge0}\mathsf E_k=\mathsf{Rep}$.
\end{proposition}
\begin{proof}
An ordinary allocation is a point-mass support selection and has one empty
support by the ample Euler identity. Nonclashing with that root gives the
rootward condition. It supplies a point in every $\mathcal F_k(C)$.
The sets $\Omega_k(C)$ are decreasing in $k$, and each $\mathcal F_k(C)$
lies in $\mathcal A(C)$, proving the inclusions.
For level zero we use the inherited
\href{../containment-rep-frac/sections/rooted_tree_transport.tex}{common-root whole-support coupling theorem}:
every full allocation is the mean of rooted actual selections. Vertex
budgets are automatic, so $\mathcal F_0(C)=\mathcal A(C)$.
If $k$ reaches the largest cube dimension, every selection in $\Omega_k(C)$
is a full integral allocation. Nonempty $\mathcal F_k(C)$ therefore supplies
an ordinary map. Each finite class has a finite largest cube dimension,
which also proves the intersection identity.
\end{proof}

\paragraph{Positivity of small-support coefficients.}
The symbol $\mathsf P_k$ already denotes chunk-peeling width in the atlas.
We use $\mathsf{Pos}_k$ for the following different hierarchy.

\begin{definition}[Dimension-filtered positivity]
\label{def:refinement-positivity}
For $k\ge0$, let
\[
 \begin{aligned}
 \mathcal Z_k(C)&=\{z\in\mathbb R^{I(C)}:M_Cz=\mathbf1,
                   \ z(v,S)\ge0\text{ whenever }|S|\le k\},\\
 \mathsf{Pos}_k&=\{C:\mathcal Z_k(C)\ne\varnothing\}.
 \end{aligned}
\]
Coefficients on larger supports are unrestricted real numbers.
These are signed full solutions, not partial systems obtained by deleting
higher cube equations.
\end{definition}

\begin{proposition}[The positivity hierarchy and its known endpoint]
\label{prop:refinement-positivity-chain}
One has
\[
 \mathsf{Frac}\subseteq\mathsf{Pos}_{k+1}\subseteq\mathsf{Pos}_k
 \subseteq\mathsf{Amp},\qquad
 \bigcap_{k\ge0}\mathsf{Pos}_k=\mathsf{Frac}.
\]
The inherited nearest-edge comparison theorem gives
$\mathsf{Pos}_2=\mathsf{Amp}$.
On classes with largest actual cube dimension at most $k$,
$\mathsf{Pos}_k=\mathsf{Frac}$.
\end{proposition}
\begin{proof}
The sign constraints increase with $k$. When $k$ reaches the largest
actual support dimension, they impose nonnegativity on every coefficient,
which proves all assertions except the known level-two equality.
For that equality we import Section~9 of
\href{../minor-atlas/research/stationary_exclusion_attack.md}{the low-order comparison proof}:
on every ample class it constructs a full signed solution nonnegative on
all supports of size zero, one and two. Its higher coefficients may be
negative. This is exactly $\mathcal Z_2(C)\ne\varnothing$.
The equality is inherited author-proved mathematics; this definition does
not supply an independent review or any higher-order positivity theorem.
\end{proof}

\begin{proposition}[Exact alternative for filtered positivity]
\label{prop:refinement-positivity-dual}
The set $\mathcal Z_k(C)$ is empty if and only if there are real prices
$y_B$, one for every actual cube, satisfying
\[
 \begin{cases}
 (M_C^{\mathsf T}y)(v,S)\ge0,& |S|\le k,\\
 (M_C^{\mathsf T}y)(v,S)=0,& |S|>k,
 \end{cases}
 \qquad \sum_{B\in\mathcal B(C)}y_B<0.
\]
\end{proposition}
\begin{proof}
Split each unrestricted higher-support variable into a difference of two
nonnegative variables. Farkas' alternative applied to the resulting
nonnegative-variable system gives inequalities on the small-support
columns and both opposite inequalities on every higher-support column.
The latter are exactly zero. The right-hand side is $\mathbf1$, yielding
the stated negative total. Conversely pairing such prices with any
$z\in\mathcal Z_k(C)$ gives a nonnegative number equal to that negative
total, a contradiction.
\end{proof}

In particular $\mathsf{Pos}_3=\mathsf{Amp}$ would give fractional
feasibility for every ample class of VC dimension at most three, since
ample shattering equals actual cube dimension. Failure of one fixed
nearest-position construction does not show absence of a different
$z\in\mathcal Z_3(C)$.

\begin{definition}[Half-integral existence; a reserve comparison]
\label{def:refinement-half-integral}
Let $\mathsf{HalfFrac}$ consist of those $C$ with an allocation in
$\mathcal A(C)$ whose coefficients belong to $\{0,\tfrac12,1\}$.
Then $\mathsf{Rep}\subseteq\mathsf{HalfFrac}\subseteq\mathsf{Frac}$.
\end{definition}
These are coefficient masses, not only edge marginals. The existing
half-integral support-face theory does not guarantee an integral point
elsewhere in the full polytope. No campaign for this reserve comparison
is initiated here, and no comparison with $\mathsf E_k$ is asserted.

\paragraph{Two whole-witness properties.}
The following notation names properties already used in the witness
comparison results; it does not assert a new equality in the original chain.
For an ample $C$ let $I(C)$ be its actual rooted columns and let
\[
 \Lambda_C=\{(N,z)\in\mathbb Z\times\mathbb Z^{I(C)}:
                         M_Cz=N\mathbf1\},\qquad
 W_C=\langle(1,\mu_r):r\text{ is a representation map of }C\rangle_{\mathbb Z},
\]
where $\mu_r(v,S)=\mathbf1_{\{r(v)=S\}}$.
The signed counts in $\Lambda_C$ use every native row and column.

\begin{definition}[Inward cube maps and full witness lattices]
\label{def:atlas-witness-properties}
Write $\mathsf{Inw}$ for the ample classes $C$ such that, for every
contained positioned cube $Q(a,S)$, there is a representation map $r$
of $C$ with $r(v)\subseteq S$ for every $v\in Q(a,S)$.
Thus every boundary edge points into that cube, simultaneously.
Write $\mathsf{Lat}=\{C\in\Amp:W_C=\Lambda_C\}$.
No nonnegativity is required of the integer coefficients generating $W_C$.
\end{definition}

\begin{proposition}[Whole-witness refinements and minor closure]
\label{prop:atlas-witness-refinements}
One has
\[
 \MaxMin\subseteq\mathsf{Inw}\subseteq\mathsf{Lat}\subseteq\Rep,
 \qquad \mathsf{CAmp}\subseteq\mathsf{Inw}.
\]
Both $\mathsf{Inw}$ and $\mathsf{Lat}$ are closed under nonempty
projections and conditionings. No strictness or comparison of
$\mathsf{Ext}$ with $\mathsf{Inw}$ or $\mathsf{Lat}$ is asserted.
\end{proposition}
\begin{proof}
The maximum-parent inward theorem
\href{../maxmin-forbidden/sections/cube_inward_representations.tex}
{\nolinkurl{thm:maximum-cube-inward}} supplies the first inclusion,
including marked sections. The cube-retaining peeling of
\cref{prop:cube-chordal-peeling} supplies the $\mathsf{CAmp}$ inclusion:
start with any map of the retained cube and reverse the peeling by
literal corner extension.

We spell out the implication to the integer lattice. The
\href{research/native_exchange_lattice.md}{integer exchange basis}
consists, for each actual cube $Q$ of directions $S$ and a spanning
tree of its edges $v w$ of direction $i$, of
\[
 e_{(v,S)}-e_{(w,S)}
 -e_{(v,S\setminus\{i\})}+e_{(w,S\setminus\{i\})}.
\]
These vectors form an integer basis of $\ker M_C$.
An inward map at $Q$ permits replacing its entire cube map while
keeping every outside column fixed: outside rows are unchanged,
and inside rows receive precisely the new cube budgets.
Inside $Q$, the distance map rooted opposite $v$ assigns $S$ to $v$
and $S\setminus\{i\}$ to $w$. Swapping those two supports changes the
array by the displayed kernel vector and retains nonnegative unit
budgets. Hence every basis vector is a difference of two complete
ordinary witnesses of $C$. One reference witness then generates the
degree direction as well, so $W_C=\Lambda_C$.

For completeness, $\Lambda_C$ always has a degree-one element.
In owner coordinates,
\[
 w_{v,T}=\sum_{S\supseteq T}z(v,S),
\]
each actual positioned cube has
the independent equation $\sum_{v\in Q}w_{v,\operatorname{dir}Q}=N$.
Choose one owner per cube and invert these Boolean upper sums
integrally. Thus $\mathsf{Lat}$ cannot have an empty ordinary witness
set, proving $\mathsf{Lat}\subseteq\Rep$.

Conditioning an inward map at a retained cube preserves inwardness.
For projection, every actual cube of $\pi_i C$ has a horizontal copy
in $C$: condition outside its directions and use strong shattering
in that ample section. Choose an inward map at that copy.
The ordinary projection selects the $i$-edge sink in every two-point
fiber. On the selected copy these are its own vertices, and their
projected supports remain inside the desired cube. This proves
projection closure of $\mathsf{Inw}$.

Finally the native conditioning and projection maps on homogeneous
signed lattices are integral surjections and send whole witnesses to
whole witnesses. Here is a direct check of the surjectivity used.
Conditioning retains exactly the owner blocks of its face, and those
blocks may be extended by arbitrary degree-$N$ owners on other cubes.
For projection, let $\alpha$ be the prescribed owner vector of a
projected cube. If it occurs in just one shore, assign $\alpha$ there.
If it occurs in both, assign $\alpha$ to both horizontal copies and
assign $\alpha$ on one endpoint shore and zero on the other for its
vertical-prism owners. The projected owner vector is then
$\alpha+\alpha-\alpha=\alpha$.
Every source cube block has total $N$, and Boolean inversion is
integral. These choices give a signed lift of every target.
The formulas for the minor maps are the
\href{../containment-campaigns/ALLOCATION_FUNCTORS.md}{native allocation minor maps}.
If $W_C=\Lambda_C$, their surjectivity therefore expresses every target
signed count as an integer combination of target whole witnesses.
This proves minor closure of $\mathsf{Lat}$.
\end{proof}

The implication to $\mathsf{Lat}$ does not decompose an arbitrary
nonnegative allocation into a probability law. Integer generation,
positive decomposition of a prescribed mean, and existence of one
ordinary map are different assertions. In particular the recent large
torsion groups attached to a \emph{selected positive support} are not
the full-witness defect $\Lambda_C/W_C$.

\begin{proposition}[Closure information for the allocation refinements]
\label{prop:atlas-refinement-minor-closure}
The families $\mathsf{BFrac}$, $\mathsf M$ and
$\mathsf{HalfFrac}$ are pc-minor closed.
Every $\mathsf{Pos}_k$ is projection closed.
No conditioning-closure assertion for $\mathsf{Pos}_k$, $k\ge3$, or
minor-closure assertion for $\mathsf E_k$, $k\ge1$, is made here.
\end{proposition}
\begin{proof}
Write $D=Q_X\setminus C$. Conditioning commutes with ambient
complementation, while
\[
 (\pi_i C)^c=\partial_i D.
\]
Fractional allocations descend through projection, conditioning and
full reduction; ordinary maps descend through projection and
conditioning. These facts prove the claims for both shores of
$\mathsf{BFrac}$ and for the ordinary/fractional requirements of
$\mathsf M$, with the allowed empty complementary output.

All those allocation maps sum source coefficients with nonnegative
integer coefficients. They preserve the grid $\tfrac12\mathbb Z$.
Nonnegative unit vertex budgets then bound each target coefficient
by one, proving pc closure of $\mathsf{HalfFrac}$.

For a signed solution $z$ on $C$, extend absent coefficients by zero.
Its projection has
\[
 z'(u,S)=z((u,0),S)+z((u,1),S),\qquad i\notin S.
\]
It satisfies all projected cube rows. If $|S|\le k$, both summands
are nonnegative for $z\in\mathcal Z_k(C)$.
Thus $z'\in\mathcal Z_k(\pi_i C)$.
Conditioning instead also uses coefficients on $S\cup\{i\}$;
at $|S|=k$ their signs are unrestricted, so this argument makes
no fixed-level conditioning claim.
\end{proof}

\paragraph{The refined existence map.}
Combining the preceding statements gives
\[
 \begin{aligned}
 \mathsf{MaxMin}&\subseteq\mathsf{Ext}\subseteq\mathsf M
 \subseteq\mathsf{Rep}\subseteq\mathsf E_2\subseteq\mathsf E_1
 \subseteq\mathsf{Frac},\\
 \mathsf{Frac}&\subseteq\mathsf{Pos}_3\subseteq\mathsf{Pos}_2
 =\mathsf{Amp}.
 \end{aligned}
\]
None of these refinements is asserted strict. A bad prescribed law or
signed construction does not separate the existential classes.
For example a full cube has the ordinary distance allocation
$r_q(v)=\{x:v(x)\ne q(x)\}$, and so lies in all the displayed families,
even though other allocations on the same cube can be nonclassical.
The independent mixed-complement, edge-coupling and cubic-positivity
workspaces own their respective existence questions; the definitions here
are a common notation contract, not three new closure theorems.

\begin{figure}[p]
\centering
{\small (a) Complementary feasibility, coupling and filtered positivity}
\par\medskip
\resizebox{.87\textwidth}{!}{%
\begin{tikzpicture}[>=Stealth,
 box/.style={draw,rounded corners,align=center,minimum width=2cm,
             minimum height=.85cm,font=\small,inner sep=4pt},
 incl/.style={->,draw=black!65,thin}]
\node[box] (amp) at (6,9) {$\Amp=\mathsf{Pos}_2$\\[-2pt]{\scriptsize\textsf{pc+}}};
\node[box] (pos) at (6,7.5) {$\mathsf{Pos}_3$\\[-2pt]{\scriptsize\textsf{p+, c?}}};
\node[box] (frac) at (6,6) {$\Frac$\\[-2pt]{\scriptsize\textsf{pc+}}};
\node[box] (bf) at (2.3,4.5) {$\mathsf{BFrac}$\\[-2pt]{\scriptsize\textsf{pc+}}};
\node[box] (e1) at (9.8,4.5) {$\mathsf E_1$\\[-2pt]{\scriptsize\textsf{p?, c?}}};
\node[box] (e2) at (9.8,3) {$\mathsf E_2$\\[-2pt]{\scriptsize\textsf{p?, c?}}};
\node[box] (half) at (13.2,3) {$\mathsf{HalfFrac}$\\[-2pt]{\scriptsize\textsf{pc+}}};
\node[box] (rep) at (9.8,1.5) {$\Rep$\\[-2pt]{\scriptsize\textsf{pc+}}};
\node[box] (mixed) at (6,0) {$\mathsf M$\\[-2pt]{\scriptsize\textsf{pc+}}};
\node[box] (ext) at (2.3,-1.5) {$\mathsf{Ext}$\\[-2pt]{\scriptsize\textsf{pc+}}};
\node[box] (mm) at (6,-3) {$\MaxMin$\\[-2pt]{\scriptsize\textsf{pc+}}};
\foreach \a/\b in {mm/ext,ext/mixed,mixed/rep,mixed/bf,bf/frac,
 rep/e2,e2/e1,e1/frac,rep/half,frac/pos,pos/amp}
 \draw[incl] (\a)--(\b);
\draw[incl] (half.north) -- (13.2,6) -- (frac.east);
\node[font=\small] at (5.9,3.1) {$\mathsf M=\Rep\cap\mathsf{BFrac}$};
\end{tikzpicture}}
\par\medskip
{\small (b) Whole-witness refinements}
\par\medskip
\resizebox{.89\textwidth}{!}{%
\begin{tikzpicture}[>=Stealth,
 box/.style={draw,rounded corners,align=center,minimum width=2cm,
             minimum height=.85cm,font=\small,inner sep=4pt},
 incl/.style={->,draw=black!65,thin}]
\node[box] (mm) at (0,0) {$\MaxMin$\\[-2pt]{\scriptsize\textsf{pc+}}};
\node[box] (camp) at (0,-1.5) {$\mathsf{CAmp}$\\[-2pt]{\scriptsize\textsf{p-, c+}}};
\node[box] (inw) at (3.5,0) {Cube-inward\\$\mathsf{Inw}$\quad{\scriptsize\textsf{pc+}}};
\node[box] (lat) at (7,0) {Full witness lattice\\$\mathsf{Lat}$\quad{\scriptsize\textsf{pc+}}};
\node[box] (rep) at (10.5,0) {$\Rep$\\[-2pt]{\scriptsize\textsf{pc+}}};
\node[box] (ext) at (7,-1.5) {$\mathsf{Ext}$\\[-2pt]{\scriptsize\textsf{pc+}}};
\foreach \a/\b in {mm/inw,camp/inw,inw/lat,lat/rep,camp/ext,ext/rep}
 \draw[incl] (\a)--(\b);
\end{tikzpicture}}
\caption{Refinements omitted from the structural overview.
Every arrow is a proved inclusion; no strictness is asserted.
The original adjacent comparisons and the added adjacent comparisons
remain undecided here. The symbol \textsf{?} means closure is not
established in this audit, not a proved failure or an assertion of
literature-wide openness. The classes in panel~(b) are defined in
\cref{def:atlas-witness-properties}; their inclusions and closure,
and the additional closure badges in panel~(a), are justified in
\cref{prop:atlas-witness-refinements,prop:atlas-refinement-minor-closure}.
No comparison of $\mathsf{Inw}$ or $\mathsf{Lat}$ with $\mathsf{Ext}$,
or of $\mathsf{HalfFrac}$ with $\mathsf E_k$, is asserted.
The $\mathsf{CAmp}\subseteq\mathsf{Ext}$ arrow follows from
\cref{prop:atlas-cube-chordal-complement}.}
\label{fig:allocation-refinements}
\end{figure}

